% =====================================================================
%  SECTIONS 9-10 — Applications, and an honest cost comparison
%  Draft 1. Every number traceable to RESULTS.md.
% =====================================================================

\section{Applications}
\label{sec:applications}

\subsection{Certified design: optimising a guaranteed floor}
\label{sec:applications:design}

Surrogate-driven design optimises a point estimate and relies on the
surrogate having been accurate at the optimum, precisely where the optimiser has been pushing hardest. Because the present construction is
differentiable and returns a guaranteed lower bound, the design objective
can be the \emph{certificate itself}.

For elasticity the objective takes a convenient form. The certified floor
on the directional Young's modulus is
$E(d) \ge 1 / (a \cdot C_{\mathrm{low}}^{-1} a)$ with
$a = \mathrm{Mandel}(d \otimes d)$, and since
$C_{\mathrm{low}} = S_{\mathrm{up}}^{-1}$ this is simply
$E(d) \ge 1/(a \cdot S_{\mathrm{up}}\, a)$. Maximising the guaranteed floor
therefore means \emph{minimising} a quadratic form in the dual Gram matrix,
requiring no inversion inside the optimisation loop.

The construction is self-penalising. A looser dual field produces a larger
$a \cdot S_{\mathrm{up}} a$ and hence a \emph{lower} guaranteed floor, so
the optimiser cannot profit from degrading the certificate. The failure
mode that makes surrogate-driven design untrustworthy is structurally
absent.

We optimise inclusion positions at a material budget held exactly fixed by
projection: after each design step a global radius scale is bisected until
the rendered volume fraction returns to target, giving projected gradient
ascent rather than a penalised objective.

Table~\ref{tab:design} reports an illustrative controlled synthetic outcome.

\begin{table}[t]
\centering
\caption{Illustrative certified design on controlled synthetic periodic
cells at exactly fixed material budget. The reported gain is in the
\emph{guaranteed} value for that specified model, certified by exact
integration.}
\label{tab:design}
\resizebox{\textwidth}{!}{%
\begin{tabular}{lrr}
\toprule
 & conductivity & elasticity \\
\midrule
objective & guaranteed $K_{\mathrm{low}}[1,1]$ & guaranteed $E_x$ floor \\
initial   & 1.9165 & 2.3701 \\
optimised & \textbf{2.9769} & \textbf{3.2346} \\
gain in the guaranteed value & \textbf{+55.3\%} & \textbf{+36.5\%} \\
final certified interval & $[2.9769, 3.0011]$ & $[3.2346, 3.3886]$ \\
final relative width & 0.6\% & 3.7\% \\
material budget drift & 0.000 & +0.00\% \\
\bottomrule
\end{tabular}}
\end{table}

Two observations apply to this controlled numerical example. First, the
budget is held exactly, so neither gain can be attributed to additional stiff phase; a soft-penalty formulation used
during development permitted a $6\%$ drift and inflated the conductivity
figure to $+57.1\%$. Second, the certificate \emph{tightened} during design
($12.6\% \to 0.6\%$ for conductivity, $9.0\% \to 3.7\%$ for elasticity),
because the trial fields continue training while the geometry moves:
raising the floor and tightening the interval are complementary rather than
competing objectives.

\subsection{Auditing conventional solvers}
\label{sec:applications:audit}

Because the interval provably contains the effective tensor for its specified
geometry, phase law, boundary conditions, tensor convention and
discretisation, a matching reference value falling outside it establishes an
inconsistency that must be resolved. At contrast
$200$ a finite-volume solution computed on the native pixel grid fell below
the certified lower bound, with
$\lambda_{\min}(K_{\mathrm{FV}} - K_{\mathrm{low}}) = -0.031$. Mesh
refinement confirmed the diagnosis: the same solver returned values inside
the interval at $h/2$ and $h/4$, approaching it from below
($+0.043$ and $+0.074$).

A rigorous two-sided certificate is therefore a convergence diagnostic for
conventional solvers, requiring no knowledge of the exact answer. We note
the honest limitation: the audit bites only when the certified interval is
tighter than the competing solver's error. In the comparison of
Section~\ref{sec:cost} the FFT estimates fell inside the intervals in every
case, and no discrepancy was detected.

\section{Cost, honestly}
\label{sec:cost}

We do not claim a speed advantage over classical homogenization solvers,
and Table~\ref{tab:cost} should be read carefully because the two methods return
different objects.

\begin{table}[t]
\centering
\caption{Certified intervals against an FFT (Moulinec--Suquet) estimate,
two-phase checkerboard. The FFT returns a point estimate with unknown
error; the certified interval provably contains the exact tensor. Exact
values $2.2361$ and $7.0711$.}
\label{tab:cost}
\resizebox{\textwidth}{!}{%
\begin{tabular}{lrrrrr}
\toprule
case & FFT estimate & FFT error & FFT time & certified interval & time \\
\midrule
contrast $5$  & 2.2320 & 0.184\% & 1.0\,s & $[2.1800, 2.2958]$ & 52\,s \\
contrast $50$ & 6.7976 & 3.867\% & 1.2\,s & $[4.5586, 10.5049]$ & 50\,s \\
\bottomrule
\end{tabular}}

\end{table}

Per query the FFT solve is roughly fifty times faster, and we make no
contrary claim. What it does not supply is an error bar: at contrast $50$
its estimate is $3.9\%$ from the analytic value, and nothing in its output
reveals this. Without the closed-form solution a practitioner has a number
and no means of bounding its error.

The regime in which the present method is preferable is many-query
certification. With $123$\,s of family training and $19$\,ms per
certificate, the amortized route overtakes repeated FFT solves at
approximately $300$ designs, and thereafter by a widening margin, while returning a guaranteed interval rather than an estimate at every query. We
present the comparison including the rows in which the FFT is preferable,
since a method that is faster only under a favourable accounting is not
faster.
